\chapter{Eigenvalues and Eigenvectors}

So far in this book a matrix has been a \emph{machine}: feed it a vector, get a
vector out. Most of the time the output points somewhere new. You push a vector
in, the matrix grabs it, turns it, stretches it, maybe flips it, and hands you
back something that points in a different direction entirely. That turning is the
whole reason matrices are interesting -- and it is also the whole reason they are
hard to think about. A single $3\times 3$ matrix can do a bewildering amount to
the vectors of $\R^3$ all at once.

This chapter is about a beautiful piece of good news. Hidden inside almost every
matrix there are a few special directions that the matrix \emph{refuses to turn}.
Along those directions the matrix does the simplest possible thing a linear map
can do: it just scales. Push a vector along one of these directions and it comes
back pointing the exact same way (or exactly backwards), only longer or shorter.
Those directions are the \emph{eigenvectors}, and the amount of stretch is the
\emph{eigenvalue}. Find them, and a complicated matrix suddenly becomes a list of
simple stretches along simple axes. That is the single most powerful idea in all
of linear algebra, and once you see it you will see it everywhere: in vibrating
bridges, in Google's PageRank, in the stability of differential equations, in the
principal components of a data set, in the rotations of computer graphics.

\section{Directions a matrix does not turn}
\label{ch23:directions}

Let us start with a picture and a single concrete matrix. Take
\[
  \mat{A}=\begin{bmatrix} 2 & 1 \\ 1 & 2 \end{bmatrix}.
\]
This is a perfectly ordinary symmetric matrix. Watch what it does to a handful of
input vectors. Apply it to $\vec{e}_1 = \colvec{1\\0}$ and you get $\colvec{2\\1}$
-- the output has swung upward, it points in a new direction. Apply it to
$\colvec{0\\1}$ and you get $\colvec{1\\2}$ -- swung the other way. For most
inputs, $\mat{A}$ turns the arrow.

But now try the diagonal direction $\vec{v}=\colvec{1\\1}$:
\[
  \mat{A}\colvec{1\\1}
  =\begin{bmatrix} 2 & 1 \\ 1 & 2 \end{bmatrix}\colvec{1\\1}
  =\colvec{2+1\\1+2}=\colvec{3\\3}=3\colvec{1\\1}.
\]
Look at that. The output $\colvec{3\\3}$ points in \emph{exactly} the same
direction as the input $\colvec{1\\1}$. The matrix did not turn this arrow at all;
it simply tripled its length. And try the anti-diagonal $\vec{w}=\colvec{1\\-1}$:
\[
  \mat{A}\colvec{1\\-1}
  =\colvec{2-1\\1-2}=\colvec{1\\-1}=1\cdot\colvec{1\\-1}.
\]
Same story -- the direction is preserved, this time with no stretch at all (the
length is unchanged because the multiplier is $1$).

\begin{figure}[H]
\centering
\begin{tikzpicture}[scale=1.2,>=Latex]
  \draw[help lines,gray!30] (-1.2,-1.2) grid (3.4,3.4);
  \draw[->] (-1.2,0)--(3.4,0) node[right]{$x$};
  \draw[->] (0,-1.2)--(0,3.4) node[above]{$y$};
  % a generic vector and its turned image
  \draw[->,vecblue,very thick] (0,0)--(1,0) node[below right]{$\vec{e}_1$};
  \draw[->,vecblue,thick,dashed] (0,0)--(2,1) node[right]{$\mat{A}\vec{e}_1$ (turned)};
  % eigenvector (1,1) -> (3,3)
  \draw[->,vecred,very thick] (0,0)--(1,1) node[above left]{$\vec{v}=\colvec{1\\1}$};
  \draw[->,vecred,thick,dashed] (0,0)--(3,3) node[above right]{$\mat{A}\vec{v}=3\vec{v}$};
  % eigenvector (1,-1) -> (1,-1)
  \draw[->,targetgreen,very thick] (0,0)--(1,-1) node[below right]{$\vec{w}=\colvec{1\\-1},\;\mat{A}\vec{w}=\vec{w}$};
\end{tikzpicture}
\caption{The matrix $\mat{A}=\protect\begin{bmatrix}2&1\\1&2\protect\end{bmatrix}$
turns most arrows (blue), but it leaves the directions $\vec{v}=\colvec{1\\1}$ and
$\vec{w}=\colvec{1\\-1}$ pointing the same way -- it only rescales them, by $3$ and
by $1$. Those are its eigenvectors.}
\label{ch23:fig-eigendirections}
\end{figure}

\begin{intuitionbox}
An eigenvector is an arrow that a matrix is content to leave \emph{on its own
line}. The matrix may make it longer, shorter, or flip it end to end, but it does
not rotate it off its line through the origin. The eigenvalue is the number that
tells you how much the matrix stretched it along that line: $3$ means ``triple
it,'' $1$ means ``leave it,'' $0.5$ means ``halve it,'' $-2$ means ``double it and
flip it around.''
\end{intuitionbox}

The word \emph{eigen} is German for ``own'' or ``characteristic'' or ``proper.''
These are the matrix's own directions, the ones intrinsic to it. Now let us write
the idea down precisely.

\begin{definition}[Eigenvalue and eigenvector]
\label{ch23:def-eigen}
Let $\mat{A}$ be an $n\times n$ matrix. A nonzero vector $\vec{x}\in\R^n$ is an
\emph{eigenvector} of $\mat{A}$ if there is a scalar $\lambda$ with
\[
  \boxed{\;\mat{A}\vec{x}=\lambda\,\vec{x}.\;}
\]
The scalar $\lambda$ is the \emph{eigenvalue} belonging to $\vec{x}$. (We allow
$\lambda$ to be complex; we always require $\vec{x}\neq\vec{0}$.)
\end{definition}

Two small but essential remarks live inside that definition.

\textbf{Why the vector must be nonzero.} If we allowed $\vec{x}=\vec{0}$, the
equation $\mat{A}\vec{0}=\lambda\vec{0}$ would be true for \emph{every} $\lambda$,
which would make the eigenvalue meaningless. The zero vector has no direction, so
it cannot tell us a direction the matrix preserves. We throw it out.

\textbf{Why the eigenvalue \emph{can} be zero.} A perfectly good eigenvalue is
$\lambda=0$. That says $\mat{A}\vec{x}=\vec{0}$ for a nonzero $\vec{x}$, i.e.
$\vec{x}$ is in the nullspace. So ``$0$ is an eigenvalue of $\mat{A}$'' is just a
new way of saying ``$\mat{A}$ is singular.'' Hold onto that; it will return.

\begin{focusbox}
$\mat{A}\vec{x}=\lambda\vec{x}$ is the most important equation in this half of the
book. Read it as: \emph{applying the matrix has the same effect as multiplying by
a single number}, but only along the special directions $\vec{x}$. The number is
$\lambda$; the direction is $\vec{x}$.
\end{focusbox}

\begin{remark}[An eigenvector is really a whole line]
\label{ch23:rem-scaling}
If $\mat{A}\vec{x}=\lambda\vec{x}$, then for any scalar $c\neq 0$ the vector
$c\vec{x}$ is an eigenvector too, with the \emph{same} eigenvalue:
$\mat{A}(c\vec{x})=c\,\mat{A}\vec{x}=c\lambda\vec{x}=\lambda(c\vec{x})$. So
eigenvectors come in lines (and, as we will see, sometimes whole planes). When we
``find the eigenvector'' we really mean ``find the line''; we usually report one
convenient representative, like $\colvec{1\\1}$ rather than $\colvec{5\\5}$.
\end{remark}

\begin{walkbox}
Before reading on, picture three transformations of the plane and ask which
directions survive untouched. (1) A pure scaling that doubles everything: which
directions are eigenvectors? (2) A reflection across the $x$-axis: which arrows
stay on their line, and what are their eigenvalues? (3) A $90^\circ$ rotation:
can \emph{any} real arrow stay on its line? Carry these on your walk; we answer
all three in this chapter.
\end{walkbox}

\section{Finding the eigenvalues: the characteristic equation}
\label{ch23:char}

Geometry told us eigenvectors exist for the matrix $\begin{bmatrix}2&1\\1&2\end{bmatrix}$
because we got lucky and guessed $\colvec{1\\1}$. We need a method that does not
rely on luck. Here is the key algebraic move.

Start from $\mat{A}\vec{x}=\lambda\vec{x}$ and drag everything to one side. The
trick is to write $\lambda\vec{x}=\lambda\mat{I}\vec{x}$ so that both sides are a
matrix times $\vec{x}$:
\[
  \mat{A}\vec{x}-\lambda\vec{x}=\vec{0}
  \quad\Longrightarrow\quad
  \mat{A}\vec{x}-\lambda\mat{I}\vec{x}=\vec{0}
  \quad\Longrightarrow\quad
  (\mat{A}-\lambda\mat{I})\,\vec{x}=\vec{0}.
\]

\begin{trapbox}
You cannot write $\mat{A}\vec{x}-\lambda\vec{x}=(\mat{A}-\lambda)\vec{x}$ and stop
there. The symbol $\mat{A}-\lambda$ is nonsense: $\mat{A}$ is a matrix and
$\lambda$ is a number, and you cannot subtract a number from a matrix. You must
insert the identity, $\lambda\vec{x}=\lambda\mat{I}\vec{x}$, turning the scalar
into the matrix $\lambda\mat{I}=\diag(\lambda,\dots,\lambda)$. \emph{Then}
$\mat{A}-\lambda\mat{I}$ is a legitimate matrix subtraction.
\end{trapbox}

Now stare at $(\mat{A}-\lambda\mat{I})\vec{x}=\vec{0}$. This is a homogeneous
system. We want a \emph{nonzero} solution $\vec{x}$. When does a square system
$\mat{M}\vec{x}=\vec{0}$ have a nonzero solution? Exactly when the matrix
$\mat{M}=\mat{A}-\lambda\mat{I}$ is \emph{singular} -- when it squashes some
nonzero vector down to $\vec{0}$, i.e. when it has a nontrivial nullspace. And we
have a perfect detector for singularity from the determinant chapters: a square
matrix is singular precisely when its determinant is zero.

\begin{formalbox}
The number $\lambda$ is an eigenvalue of $\mat{A}$ if and only if
$\mat{A}-\lambda\mat{I}$ is singular, that is, if and only if
\[
  \boxed{\;\det(\mat{A}-\lambda\mat{I})=0.\;}
\]
The left-hand side, viewed as a function of $\lambda$, is a polynomial of degree
$n$ called the \emph{characteristic polynomial} $p(\lambda)$; the equation
$p(\lambda)=0$ is the \emph{characteristic equation}. Its roots are exactly the
eigenvalues of $\mat{A}$.
\end{formalbox}

This is the engine. To find eigenvalues you do not hunt for special vectors at
all -- you compute one determinant, get a polynomial in $\lambda$, and find its
roots. The eigenvectors come afterward, one eigenvalue at a time.

\subsection{The full recipe, walked through twice}
\label{ch23:recipe}

\begin{focusbox}
\textbf{The eigenvalue recipe.}
\begin{enumerate}
  \item Form $\mat{A}-\lambda\mat{I}$ (subtract $\lambda$ from each diagonal entry).
  \item Compute $\det(\mat{A}-\lambda\mat{I})$; this is the characteristic
        polynomial $p(\lambda)$.
  \item Solve $p(\lambda)=0$ for the eigenvalues $\lambda_1,\lambda_2,\dots$
  \item For each eigenvalue $\lambda_i$, solve $(\mat{A}-\lambda_i\mat{I})\vec{x}=\vec{0}$
        to find its eigenvectors (its nullspace).
\end{enumerate}
\end{focusbox}

\begin{example}[Our running matrix, done by the book]
\label{ch23:ex-running}
Take $\mat{A}=\begin{bmatrix}2&1\\1&2\end{bmatrix}$ and rediscover its
eigenstructure without guessing.

\emph{Step 1.} Subtract $\lambda$ from the diagonal:
\[
  \mat{A}-\lambda\mat{I}=\begin{bmatrix}2-\lambda & 1\\ 1 & 2-\lambda\end{bmatrix}.
\]

\emph{Step 2.} Take the determinant of this $2\times 2$ (recall
$\det\begin{bmatrix}a&b\\c&d\end{bmatrix}=ad-bc$):
\[
  p(\lambda)=(2-\lambda)(2-\lambda)-(1)(1)=(2-\lambda)^2-1.
\]

\emph{Step 3.} Set it to zero. We could expand, but it factors nicely as a
difference of squares, $(2-\lambda)^2-1=\bigl((2-\lambda)-1\bigr)\bigl((2-\lambda)+1\bigr)=(1-\lambda)(3-\lambda)$:
\[
  p(\lambda)=(1-\lambda)(3-\lambda)=0\quad\Longrightarrow\quad
  \lambda_1=3,\ \lambda_2=1.
\]
There are the two stretch factors we found geometrically.

\emph{Step 4, eigenvalue $\lambda_1=3$.} Solve
$(\mat{A}-3\mat{I})\vec{x}=\vec{0}$:
\[
  \mat{A}-3\mat{I}=\begin{bmatrix}2-3 & 1\\ 1 & 2-3\end{bmatrix}
  =\begin{bmatrix}-1 & 1\\ 1 & -1\end{bmatrix}.
\]
The single equation both rows give is $-x_1+x_2=0$, i.e. $x_1=x_2$. So every
eigenvector for $\lambda=3$ is a multiple of $\colvec{1\\1}$. (Notice the matrix
$\begin{bmatrix}-1&1\\1&-1\end{bmatrix}$ is singular -- its two rows are
proportional -- which is exactly why $\lambda=3$ was an eigenvalue.)

\emph{Step 4, eigenvalue $\lambda_2=1$.} Solve
$(\mat{A}-1\cdot\mat{I})\vec{x}=\vec{0}$:
\[
  \mat{A}-\mat{I}=\begin{bmatrix}1 & 1\\ 1 & 1\end{bmatrix},\qquad
  x_1+x_2=0\ \Longrightarrow\ x_2=-x_1.
\]
So the eigenvectors for $\lambda=1$ are multiples of $\colvec{1\\-1}$. Both lines
match Figure~\ref{ch23:fig-eigendirections} exactly.
\end{example}

\begin{trapbox}
After you plug an eigenvalue back in, the matrix $\mat{A}-\lambda\mat{I}$ \emph{must}
come out singular -- its rows must be dependent -- or you have made an arithmetic
error. If you reduce $(\mat{A}-\lambda\mat{I})\vec{x}=\vec{0}$ and find the only
solution is $\vec{x}=\vec{0}$, something is wrong: either your $\lambda$ is not
really a root, or you slipped in the determinant. A correct eigenvalue always
leaves you a free variable.
\end{trapbox}

\begin{example}[A non-symmetric $2\times 2$]
\label{ch23:ex-nonsym}
Let $\mat{B}=\begin{bmatrix}4 & 1\\ 2 & 3\end{bmatrix}$. Then
\[
  \det(\mat{B}-\lambda\mat{I})
  =\det\begin{bmatrix}4-\lambda & 1\\ 2 & 3-\lambda\end{bmatrix}
  =(4-\lambda)(3-\lambda)-(1)(2).
\]
Expand: $(4-\lambda)(3-\lambda)=12-7\lambda+\lambda^2$, so
\[
  p(\lambda)=\lambda^2-7\lambda+12-2=\lambda^2-7\lambda+10=(\lambda-2)(\lambda-5).
\]
The eigenvalues are $\lambda=5$ and $\lambda=2$.

For $\lambda=5$: $\mat{B}-5\mat{I}=\begin{bmatrix}-1&1\\2&-2\end{bmatrix}$, giving
$-x_1+x_2=0$, so $\vec{x}=\colvec{1\\1}$.

For $\lambda=2$: $\mat{B}-2\mat{I}=\begin{bmatrix}2&1\\2&1\end{bmatrix}$, giving
$2x_1+x_2=0$, so $x_2=-2x_1$ and $\vec{x}=\colvec{1\\-2}$.

Sanity check the eigenvalue $5$:
$\mat{B}\colvec{1\\1}=\colvec{4+1\\2+3}=\colvec{5\\5}=5\colvec{1\\1}.$ Good.
\end{example}

\begin{labbox}
Compute the eigenvalues and eigenvectors of
$\mat{C}=\begin{bmatrix}3 & 0\\ 4 & 1\end{bmatrix}$ entirely by hand, following
the four steps. (Because $\mat{C}$ is triangular, you should be able to read its
eigenvalues straight off the diagonal before you even start -- see
Section~\ref{ch23:triangular}. Use the recipe anyway and confirm they agree.)
Then verify each eigenvector by direct multiplication.
\end{labbox}

\begin{strangbox}
Strang introduces eigenvalues with this very same move: ``the eigenvalues come
from $\det(\mat{A}-\lambda\mat{I})=0$, and the eigenvectors come from the
nullspace of $\mat{A}-\lambda\mat{I}$.'' If you are reading this \emph{before}
Strang, the thing to lock in is that two different ideas are stacked here: a
\emph{determinant} test finds the special numbers $\lambda$, and then a
\emph{nullspace} computation finds the special directions $\vec{x}$. Reading
\emph{after} Strang, notice we are about to push past the mechanics into what the
eigenvalues \emph{mean} -- trace, determinant, complex rotations -- which is where
the subject gets its power.
\end{strangbox}

\section{Triangular and diagonal matrices: eigenvalues for free}
\label{ch23:triangular}

Some matrices wear their eigenvalues on the outside. If $\mat{A}$ is triangular
(all zeros above, or all zeros below, the diagonal), then $\mat{A}-\lambda\mat{I}$
is also triangular, and the determinant of a triangular matrix is just the product
of its diagonal entries. So
\[
  \det(\mat{A}-\lambda\mat{I})=(a_{11}-\lambda)(a_{22}-\lambda)\cdots(a_{nn}-\lambda),
\]
and the roots are simply $\lambda=a_{11},a_{22},\dots,a_{nn}$.

\begin{proposition}
\label{ch23:prop-triangular}
The eigenvalues of a triangular (in particular, a diagonal) matrix are exactly its
diagonal entries.
\end{proposition}

\begin{example}
\label{ch23:ex-tri}
For $\mat{T}=\begin{bmatrix}5 & 8 & -2\\ 0 & 1 & 7\\ 0 & 0 & 3\end{bmatrix}$ the
eigenvalues are $5,1,3$ at a glance -- no determinant expansion needed. (The
eigenvectors still require work; only the eigenvalues are free.)
\end{example}

\begin{trapbox}
``Eigenvalues are the diagonal entries'' is true \emph{only} for triangular
matrices. For a full matrix the diagonal entries are not the eigenvalues. For
Example~\ref{ch23:ex-nonsym}, $\mat{B}$ had diagonal $4,3$ but eigenvalues $5,2$.
What \emph{is} always true is that the diagonal entries \emph{sum} to the same
thing the eigenvalues sum to -- the trace -- which is the subject of the next
section ($4+3=7=5+2$).
\end{trapbox}

\section{Trace and determinant: the eigenvalues' fingerprints}
\label{ch23:trace-det}

Here are two facts that look like coincidences in our examples and turn out to be
laws. In Example~\ref{ch23:ex-nonsym} the eigenvalues were $5$ and $2$. Their
\emph{sum} is $7$, which is $4+3$, the sum of the diagonal of $\mat{B}$ -- the
trace. Their \emph{product} is $10$, which is exactly $\det\mat{B}=4\cdot 3-1\cdot 2$.
This is no accident.

\begin{theorem}[Trace and determinant from eigenvalues]
\label{ch23:thm-trace-det}
Let $\mat{A}$ be $n\times n$ with eigenvalues $\lambda_1,\dots,\lambda_n$ (listed
with multiplicity, counting complex roots). Then
\[
  \lambda_1+\lambda_2+\cdots+\lambda_n=\tr\mat{A},
  \qquad
  \lambda_1\lambda_2\cdots\lambda_n=\det\mat{A}.
\]
The eigenvalues \emph{add up} to the trace and \emph{multiply} to the determinant.
\end{theorem}

Let us see why for the $2\times 2$ case, where the algebra is transparent and the
pattern is unmistakable. Write
$\mat{A}=\begin{bmatrix}a&b\\c&d\end{bmatrix}$. Its characteristic polynomial is
\[
  \det(\mat{A}-\lambda\mat{I})
  =(a-\lambda)(d-\lambda)-bc
  =\lambda^2-(a+d)\lambda+(ad-bc)
  =\lambda^2-(\tr\mat{A})\,\lambda+\det\mat{A}.
\]
So for any $2\times 2$ matrix the characteristic polynomial is the clean little
formula
\[
  \boxed{\;p(\lambda)=\lambda^2-(\tr\mat{A})\,\lambda+\det\mat{A}.\;}
\]
Now compare with the factored form. If the roots are $\lambda_1,\lambda_2$, then
$p(\lambda)=(\lambda-\lambda_1)(\lambda-\lambda_2)=\lambda^2-(\lambda_1+\lambda_2)\lambda+\lambda_1\lambda_2$.
Matching the coefficient of $\lambda$ gives
$\lambda_1+\lambda_2=\tr\mat{A}$, and matching the constant term gives
$\lambda_1\lambda_2=\det\mat{A}$. Done. (The general $n\times n$ proof is the same
idea: the characteristic polynomial is
$(-1)^n(\lambda-\lambda_1)\cdots(\lambda-\lambda_n)$, and expanding it, the
$\lambda^{n-1}$ coefficient is $-\sum\lambda_i=-\tr\mat{A}$ while the constant term
$p(0)=\det\mat{A}=\prod\lambda_i$.)

\begin{focusbox}
For a $2\times 2$ you can often skip the algebra entirely. You only need two
numbers: $t=\tr\mat{A}$ and $d=\det\mat{A}$. The eigenvalues are the two numbers
that add to $t$ and multiply to $d$:
\[
  \lambda=\frac{t\pm\sqrt{t^2-4d}}{2}.
\]
\end{focusbox}

\begin{example}[Reading eigenvalues off trace and determinant]
\label{ch23:ex-tracedet}
For $\mat{A}=\begin{bmatrix}2&1\\1&2\end{bmatrix}$: $t=4$, $d=4-1=3$. We need two
numbers adding to $4$ and multiplying to $3$. By inspection: $3$ and $1$. (Or use
the formula: $\lambda=\tfrac{4\pm\sqrt{16-12}}{2}=\tfrac{4\pm 2}{2}=3,1$.) Same
answer as Example~\ref{ch23:ex-running}, with no determinant of $\mat{A}-\lambda\mat{I}$
expanded.
\end{example}

\begin{aestheticbox}
There is something quietly profound here. The trace and determinant are computed
directly from the entries, with no eigenvalues in sight -- yet they are completely
determined by the eigenvalues, which are the matrix's deep, basis-independent
character. Trace and determinant are the two simplest \emph{fingerprints} of a
matrix that survive any change of coordinates. The eigenvalues are the full
fingerprint; the trace (their sum) and determinant (their product) are the first
and last things you can read off without solving the whole polynomial. This is why
they show up everywhere as quick sanity checks.
\end{aestheticbox}

\begin{labbox}
Use only trace and determinant -- do not expand a characteristic polynomial -- to
find the eigenvalues of each:
\[
  \begin{bmatrix}1&2\\2&1\end{bmatrix},\qquad
  \begin{bmatrix}0&-1\\1&0\end{bmatrix},\qquad
  \begin{bmatrix}3&0\\0&3\end{bmatrix},\qquad
  \begin{bmatrix}5&4\\0&5\end{bmatrix}.
\]
For each, the eigenvalues add to the trace and multiply to the determinant. Notice
what happens to the second one (trace $0$, determinant $1$): the discriminant
$t^2-4d=-4$ is negative. We deal with that in Section~\ref{ch23:complex}.
\end{labbox}

\section{What the eigenvalues do to the picture}
\label{ch23:geometry}

Eigenvalues are not just roots of a polynomial; each one has a geometric
personality. Once you know the eigenvalues of a $2\times 2$ matrix and its
eigenvector directions, you can sketch what the matrix does to the whole plane
without computing anything else. The reason is the trick we have been building
toward all chapter: \emph{decompose along the eigenvectors}.

Suppose $\mat{A}$ has eigenvectors $\vec{v}_1,\vec{v}_2$ that span the plane, with
eigenvalues $\lambda_1,\lambda_2$. Any vector $\vec{x}$ can be written as a
combination $\vec{x}=c_1\vec{v}_1+c_2\vec{v}_2$. Then by linearity
\[
  \mat{A}\vec{x}
  =c_1\,\mat{A}\vec{v}_1+c_2\,\mat{A}\vec{v}_2
  =c_1\lambda_1\vec{v}_1+c_2\lambda_2\vec{v}_2.
\]
The matrix acts \emph{one eigendirection at a time}: it stretches the $\vec{v}_1$
part by $\lambda_1$ and the $\vec{v}_2$ part by $\lambda_2$, independently. In
eigenvector coordinates, the most complicated thing this matrix can do is two
separate scalings. That is the whole secret, and it is the bridge to
diagonalization in the next chapter. Figure~\ref{ch23:fig-ellipse} shows the two
cases side by side: two positive eigenvalues turn the unit circle into a stretched
ellipse, while a negative eigenvalue flips one direction end-to-end.

\begin{figure}[H]
\centering
\begin{tikzpicture}[scale=0.95,>=Latex]
  % LEFT: lambda = 3, 1 (stretch along (1,1), keep (1,-1))
  \begin{scope}
    \draw[help lines,gray!25] (-2.4,-2.4) grid (2.4,2.4);
    \draw[->] (-2.5,0)--(2.6,0); \draw[->] (0,-2.5)--(0,2.6);
    % unit circle of inputs, image is ellipse along (1,1)
    \draw[vecblue,thick] (0,0) circle (1);
    \draw[targetgreen,thick,rotate=45] (0,0) ellipse (3 and 1);
    % eigendirections: stretched along (1,1) by 3, fixed along (1,-1) by 1
    \draw[->,vecred,very thick] (0,0)--({3/sqrt(2)},{3/sqrt(2)});
    \node[vecred] at (1.9,2.35) {\small $\times 3$ along $\colvec{1\\1}$};
    \draw[->,structurepurple,very thick] (0,0)--({1/sqrt(2)},{-1/sqrt(2)});
    \node[structurepurple] at (1.55,-0.95) {\small $\times 1$ along $\colvec{1\\-1}$};
    \node at (0,-2.95) {\small inputs (circle) $\mapsto$ stretched (ellipse)};
    \node at (0,2.95) {\small $\lambda_1=3,\ \lambda_2=1$};
  \end{scope}
  % RIGHT: lambda = 2, -1 (flip + stretch)
  \begin{scope}[xshift=7.2cm]
    \draw[help lines,gray!25] (-1.6,-1.6) grid (2.2,1.6);
    \draw[->] (-1.7,0)--(2.4,0); \draw[->] (0,-1.7)--(0,1.8);
    \draw[->,vecred,very thick] (0,0)--(1,0) node[below]{$\vec{v}_1$};
    \draw[->,structurepurple,very thick,dashed] (0,0)--(0,1) node[left]{$\vec{v}_2$};
    \draw[->,vecred,thick] (0,0)--(2,0) node[below right]{$\lambda_1\vec{v}_1=2\vec{v}_1$};
    \draw[->,structurepurple,thick] (0,0)--(0,-1) node[left]{$\lambda_2\vec{v}_2=-\vec{v}_2$};
    \node at (0.5,2.95) {\small $\lambda_1=2,\ \lambda_2=-1$ (one flip)};
    \node at (0.5,-2.95) {\small positive stretches, negative flips};
  \end{scope}
\end{tikzpicture}
\caption{Left: when both eigenvalues are positive, the matrix turns the unit
circle into an ellipse whose axes are the eigenvector directions, stretched by the
eigenvalues. Right: a negative eigenvalue ($\lambda_2=-1$) flips that
eigendirection end-to-end while the positive one ($\lambda_1=2$) stretches.}
\label{ch23:fig-ellipse}
\end{figure}

Here is the field guide to a single real eigenvalue $\lambda$ along its direction:

\begin{itemize}
  \item $\lambda>1$: \emph{stretch outward} along that line (the direction is
        amplified). Repeatedly applying $\mat{A}$ sends those components flying
        off to infinity.
  \item $\lambda=1$: \emph{fixed direction} -- vectors on that line do not move at
        all.
  \item $0<\lambda<1$: \emph{shrink} toward the origin along that line. Repeated
        application collapses these components to nothing.
  \item $\lambda=0$: \emph{collapse} -- the entire line is sent to the origin in
        one step. The matrix is singular; this line is the nullspace.
  \item $\lambda<0$: \emph{flip} across the origin, and stretch (if
        $\lambda<-1$) or shrink (if $-1<\lambda<0$). Each application of $\mat{A}$
        bounces that component to the other side.
\end{itemize}

\begin{intuitionbox}
Think of the eigenvalues as a verdict on the long-run fate of each eigendirection
when you apply the matrix again and again. The eigenvalue with the largest
magnitude $|\lambda|$ wins: after many applications of $\mat{A}$, almost every
vector lines up with the eigenvector of the largest $|\lambda|$, because that
component grows fastest (or shrinks slowest) relative to the others. That single
sentence is the engine behind the power method for computing eigenvectors, behind
PageRank, and behind the stability analysis of dynamical systems. We make it
precise in the diagonalization and Markov-chain chapters.
\end{intuitionbox}

\begin{applicationbox}[PageRank: the web's dominant eigenvector]
When Google ranks web pages, it builds a giant matrix $\mat{M}$ describing the
``random surfer'': start on a page, follow a random outgoing link, repeat. The
entry $M_{ij}$ is the probability of stepping from page $j$ to page $i$. The
importance of every page is the eigenvector of $\mat{M}$ belonging to the
eigenvalue $\lambda=1$ -- the one direction that the surfing process leaves
\emph{unchanged}, the steady state $\mat{M}\vec{x}=\vec{x}$. Every other
eigendirection has $|\lambda|<1$ and dies away as the surfer keeps walking, so no
matter where you start, you converge to that dominant eigenvector. PageRank is, at
heart, the sentence ``find the eigenvector for $\lambda=1$ of a matrix with a few
billion rows.'' The whole machinery of this chapter -- eigenvalue $1$ means a
direction the matrix fixes, largest $|\lambda|$ dominates the long run -- is
exactly what makes it work.
\end{applicationbox}

\begin{example}[Reading a matrix from its eigendata]
\label{ch23:ex-readoff}
Suppose someone tells you a $2\times 2$ matrix has eigenvalue $2$ along
$\colvec{1\\0}$ and eigenvalue $-1$ along $\colvec{0\\1}$. You can sketch its
action immediately: it doubles everything horizontal and flips everything vertical.
In fact it must be $\begin{bmatrix}2&0\\0&-1\end{bmatrix}$, because that diagonal
matrix has exactly those eigenvalues along the axes. Eigendata, when complete,
\emph{is} the matrix.
\end{example}

\section{Special matrices and their eigenvalues}
\label{ch23:special}

A few geometric transformations you already know have eigenvalues you can read off
from the geometry, with no algebra at all. Each is worth keeping as a mental
anchor.

\textbf{Projections.} Let $\mat{P}$ project the plane onto a line $L$ through the
origin. If $\vec{x}$ already lies on $L$, projecting does nothing:
$\mat{P}\vec{x}=\vec{x}$, eigenvalue $1$. If $\vec{x}$ is perpendicular to $L$,
projecting kills it: $\mat{P}\vec{x}=\vec{0}=0\cdot\vec{x}$, eigenvalue $0$. So
every projection has eigenvalues $1$ and $0$ -- which is exactly why projections
satisfy $\mat{P}^2=\mat{P}$ (the eigenvalues satisfy $\lambda^2=\lambda$).
Figure~\ref{ch23:fig-projection} shows the two eigendirections.

\begin{figure}[H]
\centering
\begin{tikzpicture}[scale=1.0,>=Latex]
  \draw[gray!55,thick] (-1.7,-1.7)--(1.7,1.7) node[right]{\small line $L$};
  \draw[->] (-1.8,0)--(1.8,0); \draw[->] (0,-1.8)--(0,1.8);
  % a generic vector x = (1.4, 0.2) projects onto L (y=x): proj = ((a+b)/2)(1,1)
  \coordinate (X) at (1.5,0.1);
  \coordinate (PX) at (0.8,0.8); % projection of (1.5,0.1) onto y=x is 0.8(1,1)
  \draw[->,vecblue,very thick] (0,0)--(X) node[below right]{$\vec{x}$};
  \draw[->,vecred,very thick] (0,0)--(PX) node[above left]{$\mat{P}\vec{x}$};
  \draw[gray!60,dashed] (X)--(PX);
  % on L: lambda=1
  \draw[->,targetgreen,thick] (0,0)--(-1.1,-1.1) node[below left]{on $L$: $\lambda=1$};
  % perpendicular: lambda=0 (collapses to origin)
  \draw[->,structurepurple,thick] (0,0)--(-0.85,0.85) node[above left]{$\perp L$: $\lambda=0$};
  \filldraw[structurepurple] (0,0) circle (1.6pt);
\end{tikzpicture}
\caption{Projection onto the line $L$. Vectors already on $L$ are left alone
(eigenvalue $1$); vectors perpendicular to $L$ are collapsed to the origin
(eigenvalue $0$). A generic $\vec{x}$ (blue) is dropped onto its shadow
$\mat{P}\vec{x}$ on $L$ (red) -- so it is genuinely moved and is not an eigenvector.}
\label{ch23:fig-projection}
\end{figure}

\textbf{Reflections.} Let $\mat{R}$ reflect across a line $L$. Vectors on $L$ are
untouched (eigenvalue $1$); vectors perpendicular to $L$ are sent to their
negatives (eigenvalue $-1$). A reflection always has eigenvalues $+1$ and $-1$,
which is why doing it twice is the identity: $\lambda^2=1$.
Figure~\ref{ch23:fig-reflection} shows the two eigendirections.

\begin{figure}[H]
\centering
\begin{tikzpicture}[scale=1.0,>=Latex]
  \begin{scope}
    \draw[gray!50,thick] (-1.6,-1.6)--(1.6,1.6) node[right]{\small line $L$};
    \draw[->] (-1.7,0)--(1.7,0); \draw[->] (0,-1.7)--(0,1.7);
    % reflect across y=x: (a,b) -> (b,a)
    \draw[->,vecblue,very thick] (0,0)--(1.4,0.4) node[right]{$\vec{x}$};
    \draw[->,vecred,very thick] (0,0)--(0.4,1.4) node[above]{$\mat{R}\vec{x}$};
    \draw[->,targetgreen,thick] (0,0)--(0.9,0.9) node[above right]{on $L$: $\lambda=1$};
    \draw[->,structurepurple!80!black,thick] (0,0)--(0.7,-0.7);
    \draw[->,structurepurple!80!black,thick,dashed] (0,0)--(-0.7,0.7) node[above left]{$\perp L$: $\lambda=-1$};
    \node at (0,-2.1){\small reflection across $L$};
  \end{scope}
\end{tikzpicture}
\caption{A reflection across the line $L$ fixes the direction along $L$ (eigenvalue
$+1$) and flips the perpendicular direction (eigenvalue $-1$). A generic vector
$\vec{x}$ (blue) is genuinely turned to $\mat{R}\vec{x}$ (red) -- it is not an
eigenvector.}
\label{ch23:fig-reflection}
\end{figure}

\textbf{Scalings.} The matrix $c\mat{I}$ stretches every direction by $c$. So
\emph{every} nonzero vector is an eigenvector, all with the same eigenvalue $c$.
This is the one case where eigenvectors do not pick out special lines -- because
the matrix has no preferred directions, every direction is equal.

\textbf{Shears.} The shear $\mat{S}=\begin{bmatrix}1&1\\0&1\end{bmatrix}$ slides
the plane horizontally. Its only eigenvalue is $1$ (it is triangular, both diagonal
entries $1$), and the only eigenvectors lie along the horizontal axis
$\colvec{1\\0}$ -- the layer that does not slide. A shear is the classic example
of a matrix with a \emph{repeated} eigenvalue but \emph{not enough} eigenvectors
to fill the plane (only one line, not two). We will care a great deal about that
shortfall when we ask which matrices can be diagonalized.

\begin{trapbox}
A repeated eigenvalue does \emph{not} guarantee a repeated supply of eigenvectors.
The shear has $\lambda=1$ twice but only a one-dimensional line of eigenvectors.
Contrast $\mat{I}$, which also has $\lambda=1$ twice but a whole plane of them.
``How many times $\lambda$ is a root'' (algebraic multiplicity) and ``how big the
eigenvector space is'' (geometric multiplicity) are different numbers; the second
can be smaller. This gap is the entire reason the next chapters need ``defective''
matrices and the Jordan form.
\end{trapbox}

\section{Complex eigenvalues: when a matrix is really a rotation}
\label{ch23:complex}

Return to the walking question about a $90^\circ$ rotation. Can a rotation have a
real eigenvector? Picture it: a rotation turns \emph{every} arrow by a quarter
turn. There is no arrow it leaves on its own line. So a $90^\circ$ rotation has
\emph{no real eigenvectors at all} -- and yet the characteristic equation will
still hand us two roots, because a degree-$2$ polynomial always has two roots over
the complex numbers. Those roots are complex, and they are telling us ``rotation,''
not ``stretch.''

Let us compute. The $90^\circ$ counterclockwise rotation is
\[
  \mat{Q}=\begin{bmatrix}0 & -1\\ 1 & 0\end{bmatrix}.
\]
Its trace is $0$ and its determinant is $0\cdot 0-(-1)(1)=1$, so
\[
  p(\lambda)=\lambda^2-(\tr\mat{Q})\lambda+\det\mat{Q}=\lambda^2+1=0
  \quad\Longrightarrow\quad \lambda=\pm i.
\]
The eigenvalues are the imaginary numbers $+i$ and $-i$. Notice $|\lambda|=|{\pm i}|=1$:
the rotation does not change any length, and the magnitude of a complex eigenvalue
records exactly that scaling. The \emph{angle} of the eigenvalue records the
rotation: $i=e^{i\pi/2}$ sits at angle $90^\circ$, which is precisely the angle
$\mat{Q}$ rotates by. Figure~\ref{ch23:fig-rotation} shows both pictures at once:
the turned arrows on the left, the conjugate pair $\pm i$ on the unit circle on the
right.

\begin{figure}[H]
\centering
\begin{tikzpicture}[scale=1.1,>=Latex]
  \begin{scope}
    \draw[help lines,gray!25] (-1.3,-1.3) grid (1.3,1.3);
    \draw[->] (-1.4,0)--(1.5,0) node[right]{$x$};
    \draw[->] (0,-1.4)--(0,1.5) node[above]{$y$};
    \foreach \a in {0,45,...,315}{
      \draw[->,vecblue,thick] (0,0)--(\a:1);
      \draw[->,vecred!70,thin] (\a:1)--({\a+90}:1);
    }
    \node at (0,-1.9){\small every arrow turns $90^\circ$ -- none stays on its line};
  \end{scope}
  \begin{scope}[xshift=5.6cm]
    \draw[->] (-1.4,0)--(1.5,0) node[right]{Re};
    \draw[->] (0,-1.4)--(0,1.5) node[above]{Im};
    \draw[gray!40] (0,0) circle (1);
    \filldraw[vecred] (0,1) circle (1.6pt) node[above right]{$\lambda=i$};
    \filldraw[vecred] (0,-1) circle (1.6pt) node[below right]{$\lambda=-i$};
    \draw[->,structurepurple,thick] (0,0)--(0,1);
    \coordinate (Aw) at (1,0);
    \coordinate (Ow) at (0,0);
    \coordinate (Bw) at (0,1);
    \draw pic[draw,angle radius=10mm,"$90^\circ$"] {angle=Aw--Ow--Bw};
    \node at (0,-1.9){\small eigenvalues live on the unit circle};
  \end{scope}
\end{tikzpicture}
\caption{A $90^\circ$ rotation has no real eigendirection (left): every arrow is
turned. Its eigenvalues are complex, $\pm i$, sitting on the unit circle of the
complex plane (right). The magnitude $|\lambda|=1$ says ``no stretching''; the
angle of $\lambda$ says ``rotate by $90^\circ$.''}
\label{ch23:fig-rotation}
\end{figure}

The general rotation by angle $\theta$,
\[
  \mat{Q}_\theta=\begin{bmatrix}\cos\theta & -\sin\theta\\ \sin\theta & \cos\theta\end{bmatrix},
  \qquad
  \tr=2\cos\theta,\quad \det=1,
\]
has characteristic polynomial $\lambda^2-2\cos\theta\,\lambda+1=0$, whose roots are
\[
  \lambda=\cos\theta\pm\sqrt{\cos^2\theta-1}=\cos\theta\pm i\sin\theta
  =e^{\pm i\theta}.
\]
The eigenvalues are $e^{\pm i\theta}$: unit magnitude (no stretch, $\det=1$), and
their angle is exactly the rotation angle $\theta$. This is the cleanest possible
dictionary between a matrix and a complex number.

\begin{focusbox}
\textbf{The complex eigenvalue dictionary.} A pair of complex eigenvalues
$\lambda=r e^{\pm i\theta}$ means the matrix \emph{spirals}: it rotates by the
angle $\theta$ while scaling by the radius $r=|\lambda|$ per application.
\begin{itemize}
  \item $|\lambda|=1$: pure rotation, lengths preserved.
  \item $|\lambda|<1$: inward spiral, the origin is attracting.
  \item $|\lambda|>1$: outward spiral, the origin is repelling.
\end{itemize}
Real eigenvalues stretch along fixed lines; complex eigenvalues rotate-and-scale
in invariant planes. Together they describe everything a real matrix can do.
\end{focusbox}

\begin{example}[A spiral matrix]
\label{ch23:ex-spiral}
Consider $\mat{A}=\begin{bmatrix}1 & -1\\ 1 & 1\end{bmatrix}$. Trace $2$,
determinant $1\cdot 1-(-1)(1)=2$, so
$p(\lambda)=\lambda^2-2\lambda+2=0$ and
\[
  \lambda=\frac{2\pm\sqrt{4-8}}{2}=1\pm i.
\]
Write $1\pm i=\sqrt{2}\,e^{\pm i\pi/4}$: magnitude $\sqrt 2$, angle $45^\circ$. So
$\mat{A}$ rotates by $45^\circ$ and stretches by $\sqrt2$ each time you apply it --
an outward spiral. Indeed $\mat{A}=\sqrt2\,\mat{Q}_{45^\circ}$, which you can check
directly since $\mat{Q}_{45^\circ}=\tfrac{1}{\sqrt2}\begin{bmatrix}1&-1\\1&1\end{bmatrix}$.
Repeatedly applying $\mat{A}$ to any nonzero vector traces a spiral winding out to
infinity, $45^\circ$ at a time (Figure~\ref{ch23:fig-spiral}).
\end{example}

\begin{figure}[H]
\centering
\begin{tikzpicture}[scale=0.62,>=Latex]
  \draw[help lines,gray!20] (-5,-3) grid (3,5);
  \draw[->] (-5.2,0)--(3.3,0) node[right]{$x$};
  \draw[->] (0,-3.2)--(0,5.3) node[above]{$y$};
  % orbit of (1,0) under A = sqrt2 * Q45: radius x sqrt2, angle +45 each step
  % n: point
  \coordinate (p0) at (1,0);
  \coordinate (p1) at (1,1);
  \coordinate (p2) at (0,2);
  \coordinate (p3) at (-2,2);
  \coordinate (p4) at (-4,0);
  \draw[vecblue,thick,-{Latex}] (p0)--(p1);
  \draw[vecblue,thick,-{Latex}] (p1)--(p2);
  \draw[vecblue,thick,-{Latex}] (p2)--(p3);
  \draw[vecblue,thick,-{Latex}] (p3)--(p4);
  \foreach \p/\n in {p0/0,p1/1,p2/2,p3/3,p4/4}{
    \filldraw[vecred] (\p) circle (2.4pt);
  }
  \node[below right] at (p0) {\small $\vec{x}$};
  \node[right] at (p1) {\small $\mat{A}\vec{x}$};
  \node[above right] at (p2) {\small $\mat{A}^2\vec{x}$};
  \node[above left] at (p3) {\small $\mat{A}^3\vec{x}$};
  \node[below left] at (p4) {\small $\mat{A}^4\vec{x}$};
\end{tikzpicture}
\caption{The orbit of $\vec{x}=\protect\colvec{1\\0}$ under repeated application of
the spiral matrix $\mat{A}=\protect\begin{bmatrix}1&-1\\1&1\protect\end{bmatrix}$.
Each step turns the vector $45^\circ$ counterclockwise (the angle of the complex
eigenvalue) and lengthens it by $\sqrt2$ (the magnitude $|\lambda|>1$), so the
points wind outward. Complex eigenvalues with $|\lambda|>1$ always produce an
outward spiral like this.}
\label{ch23:fig-spiral}
\end{figure}

\begin{trapbox}
When the discriminant $t^2-4d$ is negative, do not declare ``no eigenvalues.''
There are still exactly two eigenvalues; they are complex conjugates. For a
\emph{real} matrix complex eigenvalues always come in conjugate pairs
$\lambda,\bar\lambda$, because the characteristic polynomial has real coefficients
(real entries), and the complex roots of a real polynomial pair up. So a real
$2\times 2$ has either two real eigenvalues or one conjugate pair -- never exactly
one complex eigenvalue alone.
\end{trapbox}

\begin{connectionbox}[Eigenvalues, differential equations, and stability]
The verdict on $|\lambda|$ is the bridge from linear algebra to differential
equations. The system $\frac{d\vec{x}}{dt}=\mat{A}\vec{x}$ has solutions built
from $e^{\lambda t}\vec{v}$ for each eigenpair $(\lambda,\vec{v})$ -- you can check
that $\frac{d}{dt}\bigl(e^{\lambda t}\vec{v}\bigr)=\lambda e^{\lambda t}\vec{v}
=e^{\lambda t}\mat{A}\vec{v}=\mat{A}\bigl(e^{\lambda t}\vec{v}\bigr)$. The
\emph{real part} of $\lambda$ now controls fate: $\operatorname{Re}\lambda<0$ means
$e^{\lambda t}\to 0$, a stable mode that decays; $\operatorname{Re}\lambda>0$ means
blow-up; and the \emph{imaginary part} sets the frequency of oscillation. A bridge
oscillates, a circuit rings, a population cycles -- all read off from where the
eigenvalues sit in the complex plane. (We develop this fully in the differential
equations chapter; for now, notice that the same complex eigenvalues that meant
``rotation'' for $\mat{A}\vec{x}$ mean ``oscillation'' for $e^{\mat{A}t}\vec{x}$.)
\end{connectionbox}

\section{Why eigenvectors of distinct eigenvalues are independent}
\label{ch23:independence}

We keep saying ``the eigenvectors span the plane'' as if it were automatic. It is
not automatic, but there is one clean theorem that gives it to us most of the time,
and its proof is a small gem worth seeing.

\begin{theorem}[Distinct eigenvalues give independent eigenvectors]
\label{ch23:thm-indep}
If $\vec{v}_1,\dots,\vec{v}_k$ are eigenvectors of $\mat{A}$ with \emph{distinct}
eigenvalues $\lambda_1,\dots,\lambda_k$ (all different from one another), then
$\vec{v}_1,\dots,\vec{v}_k$ are linearly independent.
\end{theorem}

\begin{proof}
Take the smallest case $k=2$ first; it carries the whole idea. Suppose
$c_1\vec{v}_1+c_2\vec{v}_2=\vec{0}$. Apply $\mat{A}$ to both sides; since
$\mat{A}\vec{v}_i=\lambda_i\vec{v}_i$,
\[
  c_1\lambda_1\vec{v}_1+c_2\lambda_2\vec{v}_2=\vec{0}.
\]
Now also multiply the original relation by $\lambda_1$:
$c_1\lambda_1\vec{v}_1+c_2\lambda_1\vec{v}_2=\vec{0}$. Subtract the two:
\[
  c_2(\lambda_2-\lambda_1)\vec{v}_2=\vec{0}.
\]
Because $\vec{v}_2\neq\vec{0}$ and $\lambda_2\neq\lambda_1$, the only way out is
$c_2=0$. Back-substituting, $c_1\vec{v}_1=\vec{0}$ forces $c_1=0$ as well. So the
only combination giving $\vec{0}$ is the trivial one: $\vec{v}_1,\vec{v}_2$ are
independent. The general case is the same trick by induction: assume any $k-1$ of
them are independent; a dependence among all $k$, hit with $\mat{A}$ and combined
with a $\lambda_k$-multiple of itself, collapses to a dependence among $k-1$, which
forces all the coefficients to vanish.
\end{proof}

\begin{corollary}
\label{ch23:cor-basis}
If an $n\times n$ matrix has $n$ \emph{distinct} eigenvalues, its $n$ eigenvectors
form a basis of $\R^n$ (or $\C^n$). Every vector then decomposes into
eigendirections, and the matrix acts as $n$ independent scalings.
\end{corollary}

This is the happy case: $n$ different eigenvalues means $n$ independent
eigenvectors means a complete eigenbasis. The shear of Section~\ref{ch23:special}
fails precisely because its eigenvalues are \emph{not} distinct -- and that failure
is the cliff that the next chapter, on diagonalization, builds its guardrails
around.

\begin{strangbox}
Strang frames the punchline of this chapter as: ``eigenvectors that come from
different eigenvalues are independent, so a matrix with no repeated eigenvalues has
a full set of eigenvectors.'' If you are heading \emph{into} Strang, carry
Theorem~\ref{ch23:thm-indep} as the reason diagonalization will work; if you are
coming \emph{out} of Strang, this is the theorem you will lean on to write
$\mat{A}=\mat{X}\mat{\Lambda}\mat{X}\inv$ in the next chapter and to justify the
spectral theorem for symmetric matrices later on.
\end{strangbox}

\section{A worked $3\times 3$}
\label{ch23:threebythree}

Nothing about the recipe changes in three dimensions; the determinant is just a
little bigger. Let us do one all the way through.

\begin{example}[Eigenstructure of a $3\times 3$]
\label{ch23:ex-3x3}
Take
\[
  \mat{A}=\begin{bmatrix}2 & 0 & 0\\ 1 & 2 & 1\\ -1 & 0 & 1\end{bmatrix}.
\]
\emph{Characteristic polynomial.} Expand $\det(\mat{A}-\lambda\mat{I})$ along the
first row, which has two zeros and so is easy:
\[
  \det\begin{bmatrix}2-\lambda & 0 & 0\\ 1 & 2-\lambda & 1\\ -1 & 0 & 1-\lambda\end{bmatrix}
  =(2-\lambda)\det\begin{bmatrix}2-\lambda & 1\\ 0 & 1-\lambda\end{bmatrix}.
\]
The remaining $2\times 2$ is triangular, so its determinant is
$(2-\lambda)(1-\lambda)$. Altogether
\[
  p(\lambda)=(2-\lambda)\,(2-\lambda)(1-\lambda)=(2-\lambda)^2(1-\lambda).
\]
The eigenvalues are $\lambda=2$ (a double root) and $\lambda=1$.

\emph{Check with trace and determinant.} The trace is $2+2+1=5$, and the
eigenvalues sum to $2+2+1=5$. Good. The determinant is
$\lambda_1\lambda_2\lambda_3=2\cdot2\cdot1=4$; you can verify directly that
$\det\mat{A}=4$ by cofactor expansion. Both fingerprints match.

\emph{Eigenvectors for $\lambda=1$.} Solve $(\mat{A}-\mat{I})\vec{x}=\vec{0}$:
\[
  \mat{A}-\mat{I}=\begin{bmatrix}1 & 0 & 0\\ 1 & 1 & 1\\ -1 & 0 & 0\end{bmatrix}.
\]
Row one says $x_1=0$; row two then says $x_2+x_3=0$, so $x_2=-x_3$. Taking
$x_3=1$ gives the eigenvector $\colvec{0\\-1\\1}$.

\emph{Eigenvectors for $\lambda=2$.} Solve $(\mat{A}-2\mat{I})\vec{x}=\vec{0}$:
\[
  \mat{A}-2\mat{I}=\begin{bmatrix}0 & 0 & 0\\ 1 & 0 & 1\\ -1 & 0 & -1\end{bmatrix}.
\]
The first row is all zeros; the second and third rows both reduce to
$x_1+x_3=0$. So there are \emph{two} free variables ($x_2$ free, and $x_3$ free
with $x_1=-x_3$). The eigenvectors for $\lambda=2$ fill a whole \emph{plane},
spanned by $\colvec{0\\1\\0}$ (taking $x_2=1,x_3=0$) and $\colvec{-1\\0\\1}$
(taking $x_2=0,x_3=1$).

So even though $\lambda=2$ is a double root, here it is lucky: it supplies a
two-dimensional eigenspace, enough to pair with the $\lambda=1$ eigenvector and
fill all of $\R^3$. This matrix has a complete eigenbasis,
$\bigl\{\colvec{0\\-1\\1},\colvec{0\\1\\0},\colvec{-1\\0\\1}\bigr\}$, despite the
repeated eigenvalue. The shear had no such luck; the difference between these two
cases is exactly the algebraic-versus-geometric multiplicity gap from the trap box
in Section~\ref{ch23:special}. Figure~\ref{ch23:fig-3x3} pictures the result: a
single eigen-line for $\lambda=1$ and an entire eigen-plane for $\lambda=2$,
together filling $\R^3$.
\end{example}

\begin{figure}[H]
\centering
\begin{tikzpicture}[scale=1.0,>=Latex,line join=round]
  % simple oblique 3D axes
  \def\ax{(2.6,0)}     % x to the right
  \def\ay{(-1.1,-0.75)}% y receding
  \def\az{(0,2.6)}     % z up
  \draw[->,gray!70] (0,0)--\ax node[right]{$x_1$};
  \draw[->,gray!70] (0,0)--\ay node[below left]{$x_2$};
  \draw[->,gray!70] (0,0)--\az node[above]{$x_3$};
  % the lambda=2 eigen-plane spanned by (0,1,0) and (-1,0,1):
  % use screen-coordinate parallelogram suggesting the tilted plane
  \fill[targetgreen!18,draw=targetgreen!60]
    (-1.05,1.85)--(1.35,1.35)--(0.05,-1.85)--(-2.35,-1.35)--cycle;
  \node[targetgreen!60!black] at (1.55,1.7) {\small $\lambda=2$ plane};
  % spanning eigenvectors inside the plane
  \draw[->,vecblue,very thick] (0,0)--(-1.1,-0.75) node[below]{\small $\colvec{0\\1\\0}$};
  \draw[->,vecblue,very thick] (0,0)--(-2.6,2.6)
        node[above left]{\small $\colvec{-1\\0\\1}$};
  % the lambda=1 eigen-line through (0,-1,1): screen dir = -ay + az
  \draw[vecred,very thick] (1.1+0,-0.75-2.6) -- (-1.1,0.75+2.6);
  \draw[->,vecred,very thick] (0,0)--(1.1,2.6+0.75);
  \node[vecred] at (1.9,2.7) {\small $\lambda=1$ line};
  \filldraw (0,0) circle (1.4pt);
\end{tikzpicture}
\caption{Eigenspaces of the $3\times 3$ matrix of Example~\ref{ch23:ex-3x3}. The
double eigenvalue $\lambda=2$ owns a whole plane (green), spanned by
$\colvec{0\\1\\0}$ and $\colvec{-1\\0\\1}$; the eigenvalue $\lambda=1$ owns a single
line (red) along $\colvec{0\\-1\\1}$. The plane and the line meet only at the origin
and together span $\R^3$, so the matrix has a complete eigenbasis.}
\label{ch23:fig-3x3}
\end{figure}


\begin{labbox}
By hand, find all eigenvalues of
$\mat{A}=\begin{bmatrix}2 & 1 & 0\\ 0 & 2 & 0\\ 0 & 0 & 3\end{bmatrix}$
(triangular -- read them off), then determine the dimension of the eigenspace for
$\lambda=2$ by row-reducing $\mat{A}-2\mat{I}$. You should find $\lambda=2$ is a
double root but its eigenspace is only one-dimensional: a $3\times 3$ cousin of the
shear, a matrix that \emph{cannot} be diagonalized. Compare with
Example~\ref{ch23:ex-3x3}, where the double root was diagonalizable.
\end{labbox}

\section{Pulling it together}
\label{ch23:together}

\begin{aestheticbox}
Step back and look at what one equation, $\mat{A}\vec{x}=\lambda\vec{x}$, has
bought us. We started with a matrix as an inscrutable machine that turns every
vector differently. We found its private set of directions -- the ones it merely
scales -- and along those directions the matrix is as simple as multiplication by a
number. The eigenvalues are the matrix's deepest invariants: they add to the trace,
multiply to the determinant, are real for stretches and complex for rotations, and
forecast the long-run behavior of repeated application. From here, two roads open.
One leads to \emph{diagonalization}, $\mat{A}=\mat{X}\mat{\Lambda}\mat{X}\inv$,
where the eigenvectors become a new coordinate system in which the matrix is just a
diagonal of eigenvalues. The other leads, through symmetric matrices and the
spectral theorem, to the \emph{singular value decomposition} and the principal
components that organize modern data. Every one of those landmarks is this chapter's
single idea, applied again and again.
\end{aestheticbox}

\begin{focusbox}
\textbf{What to hold onto from this chapter.}
\begin{itemize}
  \item $\mat{A}\vec{x}=\lambda\vec{x}$ with $\vec{x}\neq\vec{0}$: eigenvectors are
        the directions a matrix does not turn; eigenvalues are how much it scales
        them.
  \item Eigenvalues are the roots of $\det(\mat{A}-\lambda\mat{I})=0$; eigenvectors
        are the nullspace of $\mat{A}-\lambda\mat{I}$.
  \item For $2\times 2$: eigenvalues add to the trace and multiply to the
        determinant; $p(\lambda)=\lambda^2-(\tr\mat{A})\lambda+\det\mat{A}$.
  \item Triangular matrices show their eigenvalues on the diagonal.
  \item $\lambda=0\Leftrightarrow$ singular; $|\lambda|$ governs long-run growth;
        complex eigenvalues $r e^{\pm i\theta}$ mean rotate-by-$\theta$,
        scale-by-$r$.
  \item Distinct eigenvalues $\Rightarrow$ independent eigenvectors $\Rightarrow$
        a complete eigenbasis (the doorway to diagonalization).
  \end{itemize}
\end{focusbox}

\section*{Exercises}
\addcontentsline{toc}{section}{Exercises}

\begin{exercise}[Warmup by hand]
\label{ch23:ex-warm}
For each matrix, find the eigenvalues and one eigenvector per eigenvalue.
\[
  \text{(a) } \begin{bmatrix}3 & 0\\ 0 & -2\end{bmatrix}\quad
  \text{(b) } \begin{bmatrix}2 & 1\\ 0 & 2\end{bmatrix}\quad
  \text{(c) } \begin{bmatrix}0 & 2\\ 2 & 0\end{bmatrix}\quad
  \text{(d) } \begin{bmatrix}1 & 4\\ 1 & 1\end{bmatrix}.
\]
\end{exercise}

\begin{exercise}[Trace and determinant only]
\label{ch23:ex-tracedet-only}
Without expanding any characteristic polynomial, find the eigenvalues of
$\begin{bmatrix}6 & 2\\ 2 & 3\end{bmatrix}$ using only that they add to the trace
and multiply to the determinant. Then verify by computing
$\det(\mat{A}-\lambda\mat{I})$ directly.
\end{exercise}

\begin{exercise}[Reflection geometry]
\label{ch23:ex-reflect}
The matrix $\mat{R}=\begin{bmatrix}0 & 1\\ 1 & 0\end{bmatrix}$ reflects the plane
across the line $y=x$. \textbf{(a)} Without computing, predict its eigenvalues and
eigenvectors from the geometry. \textbf{(b)} Confirm with the characteristic
equation. \textbf{(c)} Explain why $\mat{R}^2=\mat{I}$ follows from the eigenvalues.
\end{exercise}

\begin{exercise}[A rotation in disguise]
\label{ch23:ex-rot}
Let $\mat{A}=\begin{bmatrix}3 & -2\\ 2 & 3\end{bmatrix}$. \textbf{(a)} Show the
eigenvalues are complex and write them in the form $r e^{\pm i\theta}$. \textbf{(b)}
By what angle does $\mat{A}$ rotate, and by what factor does it scale, each time it
is applied? \textbf{(c)} What is the long-run fate of $\mat{A}^n\vec{x}$ for a
generic $\vec{x}$ -- does it spiral in or out?
\end{exercise}

\begin{exercise}[Singular means $\lambda=0$]
\label{ch23:ex-singular}
Show that the matrix $\begin{bmatrix}2 & 4\\ 1 & 2\end{bmatrix}$ has $0$ as an
eigenvalue, find the corresponding eigenvector, and explain how this eigenvector
relates to the nullspace and to the fact that $\det\mat{A}=0$.
\end{exercise}

\begin{exercise}[Eigenvalues of $\mat{A}^2$ and $\mat{A}\inv$]
\label{ch23:ex-powers}
Suppose $\mat{A}\vec{x}=\lambda\vec{x}$ with $\vec{x}\neq\vec{0}$. \textbf{(a)} Show
$\mat{A}^2\vec{x}=\lambda^2\vec{x}$, so $\lambda^2$ is an eigenvalue of
$\mat{A}^2$ with the \emph{same} eigenvector. \textbf{(b)} If $\mat{A}$ is
invertible, show $\mat{A}\inv\vec{x}=\tfrac{1}{\lambda}\vec{x}$. \textbf{(c)}
Show $(\mat{A}+c\mat{I})\vec{x}=(\lambda+c)\vec{x}$. Conclude that shifting a matrix
by $c\mat{I}$ shifts every eigenvalue by $c$ but leaves the eigenvectors alone.
\end{exercise}

\begin{exercise}[A $3\times 3$ from start to finish]
\label{ch23:ex-3full}
For $\mat{A}=\begin{bmatrix}4 & 0 & 1\\ -2 & 1 & 0\\ -2 & 0 & 1\end{bmatrix}$,
find the characteristic polynomial, all three eigenvalues, and an eigenvector for
each. Check your eigenvalues against the trace and the determinant.
\end{exercise}

\begin{exercise}[Repeated eigenvalue, two outcomes]
\label{ch23:ex-repeat}
Both $\mat{I}=\begin{bmatrix}5 & 0\\ 0 & 5\end{bmatrix}$ and
$\mat{J}=\begin{bmatrix}5 & 1\\ 0 & 5\end{bmatrix}$ have the single eigenvalue
$\lambda=5$ (doubled). For each, find the dimension of the eigenspace for
$\lambda=5$. Explain in words why one of these has a full set of eigenvectors and
the other does not, and connect this to the shear of Section~\ref{ch23:special}.
\end{exercise}

\begin{exercise}[Trace/determinant proof]
\label{ch23:ex-proof}
Prove that for any $2\times 2$ matrix $\mat{A}=\begin{bmatrix}a&b\\c&d\end{bmatrix}$,
the characteristic polynomial equals $\lambda^2-(\tr\mat{A})\lambda+\det\mat{A}$,
and deduce that the two eigenvalues satisfy $\lambda_1+\lambda_2=\tr\mat{A}$ and
$\lambda_1\lambda_2=\det\mat{A}$. Then show, as a consequence, that a real
$2\times 2$ matrix has complex (non-real) eigenvalues exactly when
$(\tr\mat{A})^2<4\det\mat{A}$.
\end{exercise}

\begin{exercise}[The Markov long run -- challenge]
\label{ch23:ex-markov}
Let $\mat{M}=\begin{bmatrix}0.8 & 0.3\\ 0.2 & 0.7\end{bmatrix}$ (each column sums
to $1$; it is a probability transition matrix). \textbf{(a)} Show $\lambda=1$ is an
eigenvalue and find its eigenvector $\vec{q}$; scale $\vec{q}$ so its entries sum to
$1$. \textbf{(b)} Find the other eigenvalue and confirm $|\lambda|<1$.
\textbf{(c)} Using the eigendecomposition idea of
Section~\ref{ch23:geometry}, argue that for \emph{any} starting probability vector
$\vec{p}_0$, the sequence $\mat{M}^n\vec{p}_0$ converges to $\vec{q}$. Why does the
eigenvalue $1$ ``win,'' and what role does the second eigenvalue play in how fast
you get there?
\end{exercise}

\begin{exercise}[Complex eigenvectors -- challenge]
\label{ch23:ex-complex-evec}
For the $90^\circ$ rotation $\mat{Q}=\begin{bmatrix}0 & -1\\ 1 & 0\end{bmatrix}$
with eigenvalue $\lambda=i$, solve $(\mat{Q}-i\mat{I})\vec{x}=\vec{0}$ over the
complex numbers to find a complex eigenvector. Then take the eigenvalue
$\lambda=-i$ and show its eigenvector is the complex conjugate of the first. (This
conjugate-pairing of complex eigenvectors mirrors the conjugate-pairing of complex
eigenvalues for real matrices.)
\end{exercise}

\begin{exercise}[Build a matrix from eigendata]
\label{ch23:ex-build}
Construct a $2\times 2$ matrix whose eigenvalues are $3$ and $-1$ with
eigenvectors $\colvec{1\\1}$ and $\colvec{1\\-1}$ respectively. \emph{Hint:} you
want $\mat{A}\colvec{1\\1}=3\colvec{1\\1}$ and $\mat{A}\colvec{1\\-1}=-\colvec{1\\-1}$;
set up and solve for the four entries, or anticipate the next chapter and build
$\mat{A}=\mat{X}\mat{\Lambda}\mat{X}\inv$ with the eigenvectors as the columns of
$\mat{X}$.
\end{exercise}
